\section{Introduction}
\label{sec:intro}

In a four-player property game, a trade between two players also changes the position of the two who are not at the table, and it changes what the counterparty can do next. Monopoly makes this concrete. A street deed earns little rent on its own, but the player who owns every deed of a colour group may build houses, and houses multiply the rent. The last missing deed of a group is therefore worth far more to the player who holds the other deeds than its board price suggests. It is also worth more to everyone else, as something to withhold. Economists describe goods of this kind as carrying identity-dependent externalities: the holder cares who acquires the object, because the acquirer's later behaviour affects the holder~\cite{jehiel1996nuclear}.

This paper asks how large language model (LLM) agents bargain over such goods when they play complete games against one another. Recent work shows that LLM negotiators can model what a counterparty wants without turning that model into better agreements~\cite{cosentino2026counterparty}, and that agents' messages in long multi-agent simulations can be checked against the simulator's state~\cite{li2026arena}. A multi-player property game adds a different demand. A competent player must reason about what an exchange enables for a rival over later turns, and must weigh that against what the exchange does for itself. We study whether the agents attend to that consequence, and what form their attention takes in offers, responses and prices.

The evidence is a fixed cohort of twenty complete four-player games played by four configured systems, which we call Sol, Gemini, Opus and Grok (Section~\ref{sec:setup}). The cohort contains 6,989 decisions, 526 negotiation threads and 920 responses to offers. Every action was chosen from a legal menu generated by a deterministic engine, and every prompt, response and event was saved. We rebuilt every offer from the request that its recipient actually received, so each response can be related to the holdings, cash and terms the responding agent saw.

Our central finding is that the agents bargain defensively over the deed that would complete a rival's group. Opening offers that would complete only the proposer's group are rejected in 138 of 187 cases (74\%), against 152 of 238 (64\%) for opening offers that complete no group. The difference is positive in each of the five seed blocks and for each of the three systems that respond under the standard response condition. When a proposer repeats the same deeds with more cash, the concession is accepted in 9 of 58 cases if the deal would complete the proposer's group, against 24 of 86 if it would complete nobody; recipients counter instead. In most refusals of offers that would complete the recipient's own group, the refusing agent's text names what the counterparty would gain (61 of 74 by two LLM-assisted readings).

The pattern has clear limits. Asking prices are higher when the deed would complete the buyer's group (median 1.63 against 1.08 times the board price), but the ordering holds in only three of five seed blocks, and comparisons of one seller pricing one deed to different buyers are mixed. Buyers do not bid consistently more for deeds that would complete their own groups. Responses also do not change with the proposer's post-trade building cash, which is the quantity the agents most often cite when they refuse. The agents' attention to rival completion is therefore visible as a category that changes their responses, not as a graded valuation of the amounts they discuss.

Completing deeds still change hands, mostly when the holder needs cash. In 14 of 22 settled cash sales of a deed that completed the buyer's group, the seller held less than \$100 before the sale, compared with 7 of 31 other cash sales. The groups assembled this way then pass through development, rent and bankruptcy. Developed groups collect almost all post-completion rent. The first bankruptcy estate of a game goes to the existing asset leader in 15 of 20 games, and the first creditor goes on to win 18 of the 20.

One configuration needs separate treatment. Opus played almost all decisions under a recovery path that followed a provider refusal of its first request. That path removed its private notes and its corrective retry, while keeping the permission to bluff. Opus's bargaining departs from the defensive pattern: as a recipient it rejects proposer-completing offers less often than other opening offers, and as a seller it prices completing deeds close to list. We report Opus separately throughout. Because Opus is also every other agent's trading partner, its condition is part of their environment, and removing it would not produce a clean comparison.

The paper makes three contributions.
\begin{itemize}
\item A population analysis of multi-player bargaining over a completion externality, built from every offer as its recipient saw it, with each contrast reported by seed block, by configuration and by response condition (Section~\ref{sec:bargaining}).
\item Evidence that the agents' strategic attention to rivals is mainly defensive and categorical. It appears in rejections, unanswered concessions, asking prices and stated reasons, but it does not track the quantities the agents cite, and completing deeds are released mainly under cash pressure (Sections~\ref{sec:bargaining} and~\ref{sec:discussion}).
\item A trajectory account that connects those bargains to development timing and rent exposure, to all 188 debt episodes including the 128 that were paid, and to the positions held just before each game's first estate transfer, with eight source-exact incidents that show the mechanisms in context (Sections~\ref{sec:consequences} and~\ref{sec:tactics}).
\end{itemize}

The study is observational. It describes what these configured systems attended to and did in these twenty games. It does not score decisions as good or bad, estimate the causal effect of a tactic, rank the underlying models, or show that any agent reasons in a particular way internally.
